\section{Konklusjon}
I denne laboratorieøvelsen ble dioders oppførsel i ulike kretser undersøkt, med fokus på klippekretser, likeretterkretser og glatting med kondensator. Resultatene viser at dioder fungerer i stor grad som forventet fra teorien, men med enkelte avvik som kan forklares ved ikke-ideelle egenskaper til dioden.
\\[1em]
For klippekretsene ble det observert at signalet ble begrenset ved omtrent 0.8V i stedet for den ideelle verdien på 0.7V, noe som skyldes diodens reelle karakteristikk. Videre viste målingene med zenerdioder at spenningsbegrensningen skjer nær de forventede nivåene, med små avvik som kan forklares ved strømavhengig spenningsfall.
\\[1em]
For likeretterkretser ble det bekreftet at enveis likeretter gir halvperiodisk utgangssignal og en DC-komponent i samsvar med teorien. Toveis likeretteren ga først avvikende resultater på grunn av jordingsfeil, men ved bruk av transformator ble resultatene i god overensstemmelse med teorien.
\\[1em]
Til slutt viste forsøket med kondensator at økt kapasitans gir redusert rippel og høyere DC-verdi, noe som bekrefter teorien om glatting i likeretterkretser.
\\[1em]
Samlet sett har forsøket gitt god forståelse for hvordan dioder kan brukes til å forme og omdanne signaler, samt hvordan praktiske forhold påvirker resultatene i forhold til idealiserte modeller.